\documentclass[border=10pt]{standalone}
\usepackage[utf8]{inputenc}
\usepackage{amsmath}
\usepackage{circuitikz}

\begin{document}
\begin{circuitikz}[european, scale=0.9, font=\small]

% Source
\draw (0, 0) to[sV, l=$V_{\mathrm{IN}}$] (0, -4) node[ground]{};
\draw (0, 0) to[R, l={RS $47\,\Omega$}] (2, 0);
\node[circ] at (2, 0) {};

% Node 2 branches up to bass row and down to treble row
\draw (2, 0) -- (2,  3);
\draw (2, 0) -- (2, -3);

% === BASS network (y = 3) ===

% R1: node 2 to node 10
\draw (2, 3) to[R, l={R1 $22\,\mathrm{k}\Omega$}] (4.5, 3);
\node[circ] at (4.5, 3) {};

% C1: node 10 to node 12, bridging the pot (routed above at y = 4.5)
\draw (4.5, 3) -- (4.5, 4.5)
    to[C, l={C1 47nF}]
    (9.5, 4.5) -- (9.5, 3);

% Bass pot: RP1 100k total. Left end (CW) = RP1HI side. mirror puts wiper below.
\draw (4.5, 3) to[pR, mirror, l={POT1 $100\,\mathrm{k}\Omega$}, n=pot1] (9.5, 3);
\node[circ] at (9.5, 3) {};
% Straight arrow inside pot body: CCW end (right) → CW end (left)
\draw[-latex, thin] (7.25, 3) -- (6.75, 3);
\node[right, font=\bfseries\small] at (pot1.wiper) {Bass};

% R2: wiper to node 50
\draw (pot1.wiper) to[R, l_={\shortstack[l]{R2\\$22\,\mathrm{k}\Omega$}}] (7, 0.5);

% R3: node 12 to output bus
\draw (9.5, 3) to[R, l={R3 $22\,\mathrm{k}\Omega$}] (14, 3);

% === TREBLE network (y = -3) ===

% R4: node 2 to node 30
\draw (2, -3) to[R, l={R4 $10\,\mathrm{k}\Omega$}] (4.5, -3);

% Treble pot: RP3 100k total. Left end (CW) = RP3HI side. Default puts wiper above.
\draw (4.5, -3) to[pR, l_={POT3 $100\,\mathrm{k}\Omega$}, n=pot3] (9.5, -3);
% Straight arrow inside pot body: CCW end (right) → CW end (left)
\draw[-latex, thin] (7.25, -3) -- (6.75, -3);
\node[right, font=\bfseries\small] at (pot3.wiper) {Treble};

% C2: wiper to node 50
\draw (pot3.wiper) to[C, l={\shortstack[l]{C2\\470pF}}] (7, -0.5);

% R5: node 32 to output bus
\draw (9.5, -3) to[R, l={R5 $10\,\mathrm{k}\Omega$}] (14, -3);

% === Summing junction (node 50) ===
% Join R2 bottom (7, 0.5) and C2 top (7, -0.5) at node 50 (7, 0)
\draw (7, 0.5) -- (7, 0) -- (7, -0.5);
\node[circ] at (7, 0) {};

% === Op-amp ===
% Output anchored at (13, 0)
\node[op amp, anchor=out] (opamp) at (13, 0) {};

% Node 50 to inverting input: horizontal then vertical
\draw (7, 0) -| (opamp.-);

% Non-inverting input to GND (short stub)
\draw (opamp.+) -- ++(0, -0.5) node[ground]{};

% === Output bus (node 3, x = 14) ===
% Short wire from op-amp output to the bus
\draw (opamp.out) -- (14, 0);

% Vertical bus connecting R3 (top), op-amp output (mid) and R5 (bottom)
\draw (14, 3) -- (14, -3);
\node[circ] at (14,  0) {};
\node[circ] at (14, -3) {};

% R99 load: node 3 to GND
\draw (14, -3) to[R, l_={\shortstack[l]{RL\\$10\,\mathrm{k}\Omega$}}] (14, -5.5) node[ground]{};

% Output terminal
\draw (14, 0) to[short, -o] (15.5, 0) node[right] {Out};

\end{circuitikz}
\end{document}
